% 1 -----------------------------------------------------------
\begin{frame}[t]{Proposed project: continuous LLM API monitoring}
\vspace{0.1cm}
\begin{center}
\begin{tikzpicture}[every node/.style={font=\small},x=1cm,y=1cm]
\useasboundingbox (-0.1,-0.05) rectangle (14.3,2.7);
\node[module,text width=2.6cm] (queries) at (1.45,0.55) {New query\\window};
\node[module,text width=2.6cm] (encoder) at (5.05,0.55) {Sentence\\encoder};
\node[accentmodule,text width=2.6cm] (mmd) at (8.65,0.55) {MMD\\distribution score};
\node[module,text width=2.6cm] (alarm) at (12.25,0.55) {Calibrated\\window decision};
\node[whitecard,text width=2.9cm,align=center,font=\small] (ref) at (8.65,2.1) {Embedded benign\\reference};
\draw[myarrow] (queries)--(encoder);
\draw[myarrow] (encoder)--(mmd);
\draw[myarrow] (mmd)--(alarm);
\draw[softarrow] (ref)--(mmd);
\end{tikzpicture}
\end{center}

\par\vspace{0.25cm}
\begin{columns}[T,onlytextwidth]
\column{0.49\textwidth}
\keyterm{Reusable artifacts}
\begin{itemize}\setlength\itemsep{0.3em}
\item Query data and baseline code.
\item Fixed BGE-small embeddings.
\item Benign-only calibration.
\end{itemize}
\column{0.48\textwidth}
\keyterm{Scope of the existing study}
\begin{itemize}\setlength\itemsep{0.3em}
\item Offline traffic windows.
\item Mixed and adaptive attacks.
\item Reference contamination.
\end{itemize}
\end{columns}
\sourceline{\refpaper. \refrepo. Published setup: query-only, offline window decisions.}
\end{frame}

% 2 -----------------------------------------------------------
\begin{frame}[t]{The research question is sustained monitoring}
\vspace{0.15cm}
\lead{How early can we detect extraction under changing normal usage, at the same false-alarm budget?}
\vspace{-0.15cm}
{\centering
\begin{tikzpicture}[x=1cm,y=1cm,every node/.style={font=\small}]
\useasboundingbox (0,-0.5) rectangle (14.3,2.4);
\node[anchor=east,align=right] at (3.05,1.7) {Changing normal usage};
\node[anchor=east,align=right] at (3.05,0.6) {Same change\\with extraction};
\fill[BlockBlue,rounded corners=2pt] (3.45,1.43) rectangle (7.3,1.97);
\fill[SoftGreen!20,rounded corners=2pt] (7.3,1.43) rectangle (13.75,1.97);
\fill[BlockBlue,rounded corners=2pt] (3.45,0.33) rectangle (7.3,0.87);
\fill[SoftGreen!20,rounded corners=2pt] (7.3,0.33) rectangle (13.75,0.87);
\node[font=\footnotesize] at (5.35,1.7) {Normal mix A};
\node[font=\footnotesize] at (10.5,1.7) {Normal mix B};
\foreach \x in {8.0,8.72,9.47,10.31,11.18,12.12,13.0}{
\draw[WarmOrange,line width=2pt] (\x,0.35)--(\x,0.85);}
\draw[gray!65,dashed] (7.3,0.15)--(7.3,2.15);
\draw[WarmOrange!80!black,dashed] (8.0,0.12)--(8.0,1.05);
\node[anchor=south,font=\footnotesize,text=WarmOrange!85!black] at (8.0,1.03) {attack onset $\tau$};
\draw[softarrow] (3.45,-0.08)--(13.85,-0.08);
\node[anchor=north,font=\footnotesize] at (8.55,-0.13) {query index};
\node[anchor=north,font=\footnotesize] at (13.7,-0.13) {$H$};
\end{tikzpicture}
\par}

\par\vspace{-0.1cm}
\begin{columns}[T,onlytextwidth]
\column{0.48\textwidth}
\keyterm{Repeated decisions}\par
{\small Window-level errors can accumulate over a monitoring run.\par}
\column{0.48\textwidth}
\keyterm{A changing reference}\par
{\small Adaptation may reduce drift alerts and also absorb extraction traffic.\par}
\end{columns}
\sourceline{Motivation: \refpaper, Limitations. The query-order schematic shows a proposed replay design, not measured outputs.}
\end{frame}

% 3 -----------------------------------------------------------
\begin{frame}[t]{Controlled replay under benign traffic changes}
\vspace{0.2cm}
\renewcommand{\arraystretch}{1.5}
{\small
\begin{tabularx}{\textwidth}{@{}L{3.0cm}Y Y@{}}
\toprule
 & \keyterm{No extraction traffic} & \keyterm{Sparse extraction traffic}\\
\midrule
\textbf{Stable normal usage} &
False alarms over a full run. &
Detection within a query budget.\\
\textbf{Changing normal usage} &
Alerts caused by benign change. &
Detection while normal usage changes.\\
\bottomrule
\end{tabularx}\par}
\vspace{0.5cm}
\keyterm{Initial pairs:}\quad
{\small\texttt{model\_leeching}\quad and\quad\texttt{Meaeq\_SST-2}}\par
\vspace{0.25cm}
Split source queries before constructing streams. Reuse each replay across methods. Measure delay in query counts.
\sourceline{\refrepo. The README illustrates id, index, and query. Production timestamps and account histories are not established by this artifact.}
\end{frame}

% 4 -----------------------------------------------------------
\begin{frame}[t]{Alarm rules and sequential baselines}
\vspace{0.1cm}
\renewcommand{\arraystretch}{1.25}
{\small
\begin{tabularx}{\textwidth}{@{}L{4.4cm}Y@{}}
\toprule
\keyterm{Comparison} & \keyterm{What it tells us}\\
\midrule
Published MMD setup &
Checks a restricted reproduction.\\
MMD + run-level threshold &
Isolates calibration for repeated decisions.\\
MMD + accumulation &
Tests evidence combined across windows.\\
KCUSUM &
An established kernel sequential detector.\\
CALM &
An established sequential-calibration approach.\\
\bottomrule
\end{tabularx}\par}
\vspace{0.35cm}
{\small\keyterm{Shared inputs:} same query pools, replays, decision times, and calibration budget.\par}
\vspace{0.15cm}
\smallnote{Keep the original threshold as a reproduction reference; calibrate eligible methods to the same run-level target for the main comparison.}
\sourceline{\refpaper;
\href{https://arxiv.org/abs/1903.01661}{Flynn \& Yoo, KCUSUM, CDC 2019};
\href{https://proceedings.mlr.press/v151/cobb22a.html}{Cobb et al., CALM, AISTATS 2022}.\\
KCUSUM and CALM use their native update and calibration schemes.}
\end{frame}

% 5 -----------------------------------------------------------
\begin{frame}[t]{Timely detection at a matched false-alarm target}
\vspace{0.2cm}
\begin{columns}[T,onlytextwidth]
\column{0.48\textwidth}
\keyterm{Benign run}
\[
\Pr_0(T\le H)
\]
{\small At least one false alarm in $H$ queries.\par}
\column{0.48\textwidth}
\keyterm{Attack run}
\[
\Pr\!\left(\tau\le T\le \min(\tau+B,H)\right)
\]
{\small Detection within $B$ queries after onset $\tau$.\par}
\end{columns}
\par\vspace{0.35cm}
{\small\keyterm{Primary comparison:} calibrate on benign replays including the declared drift family; evaluate on held-out runs. Also report delay and late or missed attacks.\par}
\vspace{0.2cm}
\begin{block}{First-alarm protocol}
{\small $T$ is the first alarm index; $T=\infty$ means no alarm.
An alarm before onset is a false alarm and fails the timely-detection criterion.}
\end{block}
\sourceline{Proposed metrics. $H$: monitoring horizon. Stable-only calibration followed by unseen drift is a separate stress test, without a matched-budget claim.}
\end{frame}

% 6 -----------------------------------------------------------
\begin{frame}[t]{Reference updates under explicit information assumptions}
\vspace{0.3cm}
\begin{columns}[c,onlytextwidth]
\column{0.58\textwidth}
\begin{tikzpicture}[x=1cm,y=1cm,every node/.style={font=\small}]
\useasboundingbox (-0.1,-0.5) rectangle (8.15,4.2);
\node[module,text width=2.0cm] (seed) at (1.2,1.8) {Trusted\\seed reference};
\node[whitecard,text width=3.55cm,align=left,minimum height=1.2cm] (fixed) at (5.9,3.45) {\textbf{Frozen}\\Keep the seed only.};
\node[orangecard,text width=3.55cm,align=left,minimum height=1.2cm] (rolling) at (5.9,1.8) {\textbf{Naive update}\\Admit recent queries.};
\node[greencard,text width=3.55cm,align=left,minimum height=1.2cm] (trusted) at (5.9,0.15) {\textbf{Verified update}\\Admit confirmed benign.};
\draw[softarrow] (seed.east)--++(0.6,0)|-(fixed.west);
\draw[softarrow] (seed.east)--(rolling.west);
\draw[softarrow] (seed.east)--++(0.6,0)|-(trusted.west);
\end{tikzpicture}

\column{0.39\textwidth}
\keyterm{Proposed hypothesis}\par
{\small Updates may reduce alerts from benign change, while unverified updates may hide persistent extraction.\par}
\vspace{0.3cm}
\keyterm{Information requirement}\par
{\small Trusted updates need verified benign queries. Hidden test labels define an oracle diagnostic only.\par}
\end{columns}
\sourceline{Reference contamination is evaluated in \refpaper. Proposed comparison: full-run update policies under drift, with a shared verification budget.}
\end{frame}

% 7 -----------------------------------------------------------
\begin{frame}[t]{Evidence required for a research contribution}
\vspace{0.2cm}
\renewcommand{\arraystretch}{1.4}
{\small
\begin{tabularx}{\textwidth}{@{}L{7.0cm}Y@{}}
\toprule
\keyterm{Evidence from the study} & \keyterm{Next step}\\
\midrule
A failure survives matched calibration and existing sequential baselines. &
Investigate a targeted monitoring method.\\
Reference updates introduce a further failure under explicit assumptions. &
Study the update rule and information needs.\\
Simple calibration or existing methods resolve the difficulty. &
Report it; narrow or revise the question.\\
\bottomrule
\end{tabularx}\par}
\vspace{0.5cm}
Check whether the pattern survives a second data pair and a new source split.
\sourceline{Decision criteria for a proposed study. No new performance results or novelty claims are assumed here.}
\end{frame}

% 8 -----------------------------------------------------------
\begin{frame}[t]{Study scope and expected outputs}
\vspace{0.2cm}
\lead{Query-only monitoring under declared benign changes.}
\begin{columns}[T,onlytextwidth]
\column{0.48\textwidth}
\keyterm{Experimental scope}
\begin{itemize}\setlength\itemsep{0.55em}
\item Controlled query-only replay.
\item Fixed query representations.
\item Declared drift and verification.
\end{itemize}
\column{0.48\textwidth}
\keyterm{Expected outputs}
\begin{itemize}\setlength\itemsep{0.55em}
\item Reproducible replay configurations.
\item False alarms and timely detection.
\item Evidence for a focused method.
\end{itemize}
\end{columns}
\par\vspace{0.25cm}
{\small Residual failures determine whether a new monitoring mechanism is needed.\par}
\sourceline{Proposed study. No new detector performance or deployment guarantee is assumed.}
\end{frame}
